% from ega4-16-fr.tex
\oldpage[IV-4]{38}
Si $\mathscr{F}$ est un $\mathcal{O}_X$-Module quasi-cohérent, il en est de
même de $\mathscr{P}_{X/S}^n(\mathscr{F})$ pour l'une ou l'autre de ses
structures de $\mathcal{O}_X$-Module. Si de plus $\mathscr{F}$ est de type
fini (resp. de présentation finie) et $f:X\to S$ localement de type fini
(resp. localement de présentation finie),
$\mathscr{P}_{X/S}^n(\mathscr{F})$ est (pour l'une ou l'autre de ses
structures de $\mathcal{O}_X$-Module) de type fini (resp. de présentation
finie), comme il résulte de (16.3.9) et de (16.4.22).
\end{env}

\begin{env}[16.7.5]
\label{IV.16.7.5}
La définition (16.7.2.1) entraîne l'existence d'un homomorphisme de faisceaux
de groupes commutatifs
\[
  d_{X/S,\mathscr{F}}^n:\mathscr{F}
  \longrightarrow\mathscr{P}_{X/S}^n(\mathscr{F})
  \qquad\text{(aussi noté $d_{X/S}^n$)}
  \tag{16.7.5.1}
  \label{IV.16.7.5.1}
\]
tel qu'avec les notations de (16.7.4), on ait
\[
  d_{X/S,\mathscr{F}}^n(t)=1\otimes t
  \tag{16.7.5.2}
  \label{IV.16.7.5.2}
\]
et par suite, en vertu de (16.7.4.1)
\[
  d_{X/S,\mathscr{F}}^n(at)
  =(1\otimes t)a=(d_{X/S,\mathscr{F}}^n(t)).a
  \tag{16.7.5.3}
  \label{IV.16.7.5.3}
\]
\[
  d_{X/S,\mathscr{F}}^n(at)
  =(d_{X/S}^n(a)).(1\otimes t)
  =(d_{X/S}^n(a))(d_{X/S,\mathscr{F}}^n(t)).
  \tag{16.7.5.4}
  \label{IV.16.7.5.4}
\]
Il est donc $\mathcal{O}_X$-linéaire pour la structure de
$\mathcal{O}_X$-Module à droite sur $\mathscr{P}_{X/S}^n(\mathscr{F})$, et
\emph{semi-linéaire} (relativement à l'automorphisme $\sigma$ (16.3.4)) pour
la structure de $\mathcal{O}_X$-Module à gauche.
\end{env}

\begin{proposition}[16.7.6]
\label{IV.16.7.6}
Le $\mathcal{O}_X$-Module à gauche $\mathscr{P}_{X/S}^n(\mathscr{F})$ est
engendré par l'image de $\mathscr{F}$ par l'homomorphisme canonique
$d_{X/S,\mathscr{F}}^n$.
\end{proposition}

Cela résulte aussitôt de (16.7.5.3) et du cas particulier correspondant à
$\mathscr{F}=\mathcal{O}_X$ (16.3.8).

\begin{env}[16.7.7]
\label{IV.16.7.7}
Les homomorphismes canoniques de faisceaux d'anneaux
\[
  \varphi_{nm}:\mathscr{P}_{X/S}^m\longrightarrow\mathscr{P}_{X/S}^n
\]
pour $n\leq m$ (16.1.2) définissent, en vertu de (16.7.2.1) des
homomorphismes canoniques
\[
  \mathscr{P}_{X/S}^m(\mathscr{F})
  \longrightarrow\mathscr{P}_{X/S}^n(\mathscr{F})qquad(n\leq m)
\]
qui sont des homomorphismes de $\mathcal{O}_X$-Bimodules en vertu de (16.1.6)
et (16.7.4.1); en outre on a des diagrammes commutatifs
\[
  \xymatrix{
    \mathscr{P}_{X/S}^m(\mathscr{F})\ar[rr]
      &&\mathscr{P}_{X/S}^n(\mathscr{F})\\
    &\mathscr{F}\ar[ul]^-{d_{X/S,\mathscr{F}}^m}
      \ar[ur]_-{d_{X/S,\mathscr{F}}^n}
  }
\]
On a donc ainsi un système projectif de $\mathcal{O}_X$-Bimodules
$(\mathscr{P}_{X/S}^n(\mathscr{F}))$, et l'on pose
\[
  \mathscr{P}_{X/S}^{\infty}(\mathscr{F})
  =\varprojlim\mathscr{P}_{X/S}^n(\mathscr{F}).
  \tag{16.7.7.1}
  \label{IV.16.7.7.1}
\]
En outre, ce qui précède montre que les homomorphismes (16.7.5.1) forment un
système projectif d'homomorphismes, et définissent donc un homomorphisme
canonique
\[
  d_{X/S,\mathscr{F}}^{\infty}:\mathscr{F}
  \longrightarrow\mathscr{P}_{X/S}^{\infty}(\mathscr{F}).
  \tag{16.7.7.2}
  \label{IV.16.7.7.2}
\]
\end{env}

